\documentclass{article}
\usepackage[T1]{fontenc}
\usepackage{lmodern}
\usepackage{graphicx}
\usepackage{amsmath,amssymb}
\usepackage[a4paper,margin=2cm]{geometry}
\usepackage[absolute,overlay]{textpos}
\setlength{\TPHorizModule}{1mm}
\setlength{\TPVertModule}{1mm}
\usepackage[indonesian]{babel}
\usepackage{hyperref}
\setlength{\emergencystretch}{2em}
% unit: Halaman 32

\begin{document}

% === Halaman 32 (llm) ===

\begin{itemize}
  \item \textit{Cipher} ini berhasil dipecahkan oleh Babbage dan Kasiski pada pertengahan Abad 19 (akan dijelaskan pada materi selanjutnya).
  \item Kasiski menguraikan langkah-Langkah untuk menemukan panjang kunci (bukan huruf-huruf kunci kunci ).
  \item \textit{Vig\grave{e}nere Cipher} digunakan oleh Tentara Konfederasi (\textit{Confederate Army}) pada Perang Sipil Amerika (\textit{American Civil war}).
  \item Perang Sipil terjadi setelah \textit{Vig\grave{e}nere Cipher} berhasil dipecahkan.
\end{itemize}

\end{document}
